\chapter{Decomposition does not change the answer}

A linear discrete problem $Au=f$ is cut into pieces.  Each piece is solved
exactly, given data on the cut.  The question of this chapter: when the pieces
have been made to agree, is the assembled field the solution of the undivided
problem?

For a direct interface solve and for Dirichlet--Neumann the answer is yes, with
no further hypothesis.  For restricted additive Schwarz it is yes under one
hypothesis, and false without it.

\section{T3a: a direct Schur interface solve}

\paragraph{In plain words.}
Split the unknowns into those inside the pieces and those on the interface.
Substructuring solves a smaller system for the interface unknowns alone, then
each piece solves for its interior with that interface data.  The result is
the undivided solution: nothing is lost and nothing is added by the cut.

\begin{theorem}[T3a, direct Schur solve]\label{thm:T3a}
  \lean{Atlas.schur_iff}
  \leanok
  Let $V_I,V_\Gamma$ be vector spaces over a field, and
  $A_{II}:V_I\to V_I$, $A_{I\Gamma}:V_\Gamma\to V_I$,
  $A_{\Gamma I}:V_I\to V_\Gamma$, $A_{\Gamma\Gamma}:V_\Gamma\to V_\Gamma$ linear,
  with $A_{II}$ invertible.  For all $f_I,u_I\in V_I$ and
  $f_\Gamma,u_\Gamma\in V_\Gamma$,
  \[
    \begin{cases}A_{II}u_I+A_{I\Gamma}u_\Gamma=f_I\\
    A_{\Gamma I}u_I+A_{\Gamma\Gamma}u_\Gamma=f_\Gamma\end{cases}
    \iff
    \begin{cases}u_I=A_{II}^{-1}(f_I-A_{I\Gamma}u_\Gamma)\\
    \bigl(A_{\Gamma\Gamma}-A_{\Gamma I}A_{II}^{-1}A_{I\Gamma}\bigr)u_\Gamma
      =f_\Gamma-A_{\Gamma I}A_{II}^{-1}f_I .\end{cases}
  \]
\end{theorem}

\begin{proof}
  \leanok
  The first equation on the left is equivalent to the first on the right
  because $A_{II}$ is invertible.  Substituting it into the second equation on
  the left gives the second on the right, and conversely.
\end{proof}

\section{T3b: Dirichlet--Neumann}

\paragraph{In plain words.}
Two pieces meet along an interface and do not overlap.  One sweep: piece 1 is
solved with the interface values given; the flux it then sends through the
interface is handed to piece 2, which returns its own interface values; the
interface values move toward those by a relaxation factor $\theta$.  If a sweep
returns the interface values it was given, the two pieces' solutions side by
side solve the undivided problem.  The solves must be exact.  Convergence of
the iteration is not assumed.

\begin{definition}[Two pieces in block form]\label{def:two-pieces}
  \lean{Atlas.TwoPieces, Atlas.TwoPieces.Solves, Atlas.TwoPieces.dnStep}
  \leanok
  The unknowns are $(u_1,u_2,\lambda)\in V_1\times V_2\times V_\Gamma$.  The
  undivided system is
  \begin{align*}
    A_{11}u_1+A_{1\Gamma}\lambda&=f_1,\\
    A_{22}u_2+A_{2\Gamma}\lambda&=f_2,\\
    A_{\Gamma1}u_1+A_{\Gamma2}u_2
      +\bigl(A^{(1)}_{\Gamma\Gamma}+A^{(2)}_{\Gamma\Gamma}\bigr)\lambda
      &=f^{(1)}_\Gamma+f^{(2)}_\Gamma ,
  \end{align*}
  where each piece contributes its own share of the interface block and of the
  interface right-hand side.  The flux piece 1 sends is
  $q(u_1,\lambda)=A_{\Gamma1}u_1+A^{(1)}_{\Gamma\Gamma}\lambda-f^{(1)}_\Gamma$.
  A Dirichlet solve $D$ is exact if $A_{11}D(\lambda)+A_{1\Gamma}\lambda=f_1$
  for all $\lambda$.  A Neumann solve $N$, returning a pair
  $(u_2,\tilde\lambda)$, is exact if
  $A_{22}u_2+A_{2\Gamma}\tilde\lambda=f_2$ and
  $A_{\Gamma2}u_2+A^{(2)}_{\Gamma\Gamma}\tilde\lambda=f^{(2)}_\Gamma-q$ for all
  $q$.  One sweep is
  $\lambda\mapsto\lambda+\theta\,(\tilde\lambda-\lambda)$ with
  $(u_2,\tilde\lambda)=N\bigl(q(D(\lambda),\lambda)\bigr)$.
\end{definition}

\begin{theorem}[T3b, Dirichlet--Neumann, block form]\label{thm:T3b}
  \lean{Atlas.TwoPieces.dn_fixedPoint_solves}
  \leanok
  \uses{def:two-pieces}
  With exact solves and $\theta\ne0$, if $\lambda$ is a fixed point of the
  sweep then $\bigl(D(\lambda),u_2,\lambda\bigr)$ solves the undivided system,
  where $u_2$ is the interior the Neumann solve returned.
\end{theorem}

\begin{proof}
  \uses{def:two-pieces}
  \leanok
  $\theta\ne0$ gives $\tilde\lambda=\lambda$.  The two interior equations are
  the exactness of the solves.  The Neumann solve's interface equation,
  $A_{\Gamma2}u_2+A^{(2)}_{\Gamma\Gamma}\lambda=f^{(2)}_\Gamma-q$, with $q$
  written out, is the interface equation of the undivided system.
\end{proof}

\begin{theorem}[T3b, converse]\label{thm:T3b-converse}
  \lean{Atlas.TwoPieces.dn_solution_isFixedPt}
  \leanok
  \uses{def:two-pieces}
  If each piece's equations determine its solution and the solves return it,
  then for every solution $(u_1,u_2,\lambda)$ of the undivided system and every
  $\theta$, $\lambda$ is a fixed point of the sweep.
\end{theorem}

\begin{proof}
  \uses{def:two-pieces}
  \leanok
  The Dirichlet solve returns $u_1$.  The pair $(u_2,\lambda)$ satisfies the
  Neumann piece's equations with the flux $q(u_1,\lambda)$, so the Neumann solve
  returns it, and the update is zero.
\end{proof}

\begin{theorem}[T3b, the face form the workbench implements]\label{thm:T3b-face}
  \lean{Atlas.FacePair.solves_of_fixedPoint}
  \leanok
  Two sets of finite-volume cells, $D$ and $N$, meet along faces.  $K_D,K_N$ are
  each set's operator with the cut faces left out; $P_D,P_N$ read the cell value
  beside each cut face; $Q_D,Q_N$ put a per-face flow into those cells'
  equations; $g_D$ is the $D$ side's half-cell conductance, $r_D=g_D^{-1}$ and
  $r_N$ the two half-cell resistances, and $h=(r_D+r_N)^{-1}$ the undivided
  grid's face conductance (the harmonic mean).  Suppose
  \begin{enumerate}
    \item $K_Du_D+Q_Dg_DP_Du_D=b_D+Q_Dg_D\lambda$ (the $D$ set solved with face
      values $\lambda$);
    \item $q=g_D(P_Du_D-\lambda)$ (the flow it sends);
    \item $K_Nu_N=b_N+Q_Nq$ (the $N$ set solved with that flow);
    \item $P_Nu_N+r_Nq=\lambda$ (the $N$ set's face values reproduce
      $\lambda$).
  \end{enumerate}
  Then $q=h\,(P_Du_D-P_Nu_N)$, and $(u_D,u_N)$ solves the undivided system
  \[
    K_Du_D+Q_Dh(P_Du_D-P_Nu_N)=b_D,\qquad
    K_Nu_N-Q_Nh(P_Du_D-P_Nu_N)=b_N .
  \]
\end{theorem}

\begin{proof}
  \leanok
  From 2, $r_Dq=P_Du_D-\lambda$; from 4, $r_Nq=\lambda-P_Nu_N$.  Adding,
  $(r_D+r_N)q=P_Du_D-P_Nu_N$, so $q=h(P_Du_D-P_Nu_N)$.  Then 1 reads
  $K_Du_D+Q_Dq=b_D$ and 3 reads $K_Nu_N-Q_Nq=b_N$.
\end{proof}

\section{T3c: restricted additive Schwarz}

\paragraph{In plain words.}
The domain is covered by overlapping windows.  One sweep solves every window
exactly, with the current iterate held outside it, and blends the windows'
solutions with weights that sum to one.

\begin{itemize}
  \item The undivided solution is always a fixed point of the sweep.
  \item If, at $u$, every window's own solution agrees with $u$ on the whole
    window, then $u$ is the undivided solution.
  \item If only the \emph{blend} is unchanged, more is needed.  The sweep is
    $u\mapsto u+M(f-Au)$, with $M$ the blended sum of the window solves, and a
    fixed point is the undivided solution exactly when $M$ sends no non-zero
    residual to zero.  That holds whenever the sweep converges from every
    start.
  \item Without that hypothesis the claim is false
    (Theorem~\ref{thm:T3c-counterexample}).
\end{itemize}

\begin{definition}[A Schwarz sweep]\label{def:schwarz}
  \lean{Atlas.Schwarz, Atlas.Schwarz.window, Atlas.Schwarz.sweep, Atlas.Schwarz.precond}
  \leanok
  $A:V\to V$ is linear.  For each window $i$ of a finite family:
  $R_i:V\to W_i$ reads the window's unknowns; $E_i:W_i\to V$ pads with zeros;
  $E^\chi_i:W_i\to V$ multiplies by the window's weights, then pads;
  $B_i:W_i\to W_i$ is the inverse of the window's operator $R_iAE_i$.  The
  weights are a partition of unity: $\sum_iE^\chi_iR_i=I$.  The window's
  solution, the sweep and the preconditioner are
  \[
    w_i(f,u)=B_i\bigl(R_if-R_iA(u-E_iR_iu)\bigr),\qquad
    G_f(u)=\sum_iE^\chi_i\,w_i(f,u),\qquad
    M=\sum_iE^\chi_iB_iR_i .
  \]
\end{definition}

\begin{lemma}[The sweep in residual form]\label{lem:T3c-sweep}
  \lean{Atlas.Schwarz.sweep_eq}
  \leanok
  \uses{def:schwarz}
  $G_f(u)=u+M(f-Au)$.
\end{lemma}

\begin{proof}
  \uses{def:schwarz}
  \leanok
  $w_i(f,u)=B_iR_i(f-Au)+B_i(R_iAE_i)R_iu=B_iR_i(f-Au)+R_iu$.  Blend, and use
  the partition of unity.
\end{proof}

\begin{theorem}[T3c, consistency]\label{thm:T3c-consistency}
  \lean{Atlas.Schwarz.solution_isFixedPt}
  \leanok
  \uses{lem:T3c-sweep}
  If $Au=f$ then $G_f(u)=u$.
\end{theorem}

\begin{proof}
  \uses{lem:T3c-sweep}
  \leanok
  $M0=0$.
\end{proof}

\begin{lemma}[What a fixed point is]\label{lem:T3c-iff}
  \lean{Atlas.Schwarz.isFixedPt_iff}
  \leanok
  \uses{lem:T3c-sweep}
  $G_f(u)=u\iff M(f-Au)=0$.
\end{lemma}

\begin{proof}
  \uses{lem:T3c-sweep}
  \leanok
  Lemma~\ref{lem:T3c-sweep}.
\end{proof}

\begin{theorem}[T3c, agreement]\label{thm:T3c-agree}
  \lean{Atlas.Schwarz.solves_of_windows_agree}
  \leanok
  \uses{def:schwarz}
  If $w_i(f,u)=R_iu$ for every window $i$, then $Au=f$.
\end{theorem}

\begin{proof}
  \uses{def:schwarz}
  \leanok
  As in Lemma~\ref{lem:T3c-sweep}, $w_i(f,u)=R_iu+B_iR_i(f-Au)$, so
  $B_iR_i(f-Au)=0$ and, $B_i$ being invertible, $R_i(f-Au)=0$ for every $i$.
  By the partition of unity $f-Au=\sum_iE^\chi_iR_i(f-Au)=0$.
\end{proof}

\begin{theorem}[T3c, a stationary blend]\label{thm:T3c-precond}
  \lean{Atlas.Schwarz.fixedPt_solves_of_precond}
  \leanok
  \uses{lem:T3c-iff}
  If $Mr=0$ implies $r=0$, then every fixed point of $G_f$ solves $Au=f$.
\end{theorem}

\begin{proof}
  \uses{lem:T3c-iff}
  \leanok
  Lemma~\ref{lem:T3c-iff}.
\end{proof}

\begin{corollary}[The same hypothesis, read off the sweep]\label{cor:T3c-homogeneous}
  \lean{Atlas.Schwarz.fixedPt_solves}
  \leanok
  \uses{thm:T3c-precond, lem:T3c-sweep}
  If $A$ is onto and the sweep with zero right-hand side has no fixed point but
  $0$, then for every $f$ every fixed point of $G_f$ solves $Au=f$.
\end{corollary}

\begin{proof}
  \uses{thm:T3c-precond, lem:T3c-sweep}
  \leanok
  Let $Mr=0$ and write $r=Ae$.  Then $G_0(-e)=-e+M(Ae)=-e$, so $e=0$ and
  $r=0$.
\end{proof}

\begin{theorem}[T3c, a convergent sweep has one fixed point]\label{thm:T3c-unique}
  \lean{Atlas.Schwarz.fixedPt_eq_solution}
  \leanok
  \uses{lem:T3c-sweep}
  Let $V$ be a real normed space, and suppose $G_0^{\,k}(e)\to0$ for every
  $e\in V$.  If $Au^\star=f$ and $G_f(u)=u$, then $u=u^\star$.
\end{theorem}

\begin{proof}
  \uses{lem:T3c-sweep}
  \leanok
  $G_f(u)-G_f(u^\star)=G_0(u-u^\star)$, so $e=u-u^\star$ is a fixed point of
  $G_0$.  Then $e=G_0^{\,k}(e)\to0$.
\end{proof}

\begin{theorem}[T3c, a convergent sweep converges to the solution]\label{thm:T3c-tendsto}
  \lean{Atlas.Schwarz.tendsto_solution}
  \leanok
  \uses{lem:T3c-sweep, thm:T3c-consistency}
  Under the same hypothesis, if $Au^\star=f$ then $G_f^{\,k}(u_0)\to u^\star$
  for every $u_0\in V$.
\end{theorem}

\begin{proof}
  \uses{lem:T3c-sweep, thm:T3c-consistency}
  \leanok
  $G_f^{\,k}(u_0)-u^\star=G_0^{\,k}(u_0-u^\star)$.
\end{proof}

\begin{theorem}[T3c, the counterexample]\label{thm:T3c-counterexample}
  \lean{Atlas.Schwarz.exists_spurious_fixedPt}
  \leanok
  \uses{def:schwarz}
  There is a Schwarz sweep over $\Q$ with exact window solves, a partition of
  unity and an invertible $A$, and a pair $f,u$ with $G_f(u)=u$ and $Au\ne f$.
\end{theorem}

\begin{proof}
  \uses{def:schwarz}
  \leanok
  Three unknowns, windows $\{1,2\}$ and $\{2,3\}$, the shared unknown given
  entirely to the first window:
  \[
    A=\begin{pmatrix}1&1&0\\1&0&1\\0&1&1\end{pmatrix},\qquad f=0,\qquad
    u=(1,-1,-1).
  \]
  $\det A=-2$ and both window matrices have determinant $-1$.  Window 1 returns
  $(1,-1)$, which is $u$ on $\{1,2\}$.  Window 2 returns $(1,-1)$ on $\{2,3\}$:
  it disagrees with $u$ at the shared unknown, and the blend discards that
  value.  So $G_0(u)=u$, while $Au=(0,0,-2)$.
\end{proof}

\begin{remark}[Positive definite matrices]
  What fails in the counterexample is that $A$ has a zero on its diagonal at
  the shared unknown: the problem restricted to the overlap is singular.  For
  \emph{two} windows and a symmetric positive definite $A$ this cannot happen,
  and every fixed point is the solution (a paper proof, not yet formalised).
  With \emph{three} windows positive definiteness is not enough: the project's
  script \texttt{scripts/lean\_t3\_counterexample.py} exhibits, in exact
  rational arithmetic, a symmetric positive definite $6\times6$ matrix and three
  windows with a spurious fixed point.  On that example the sweep does not
  converge, so Theorem~\ref{thm:T3c-unique} is not contradicted.  For
  M-matrices the sweep always converges (Frommer and Szyld, \emph{SIAM J.
  Numer. Anal.} 39, 2001).
\end{remark}

\section{T3d: what stopping on the update leaves}

\paragraph{In plain words.}
A sweep is stopped when its update is small.  The update is not the error.  If
the zero-right-hand-side sweep shrinks a norm by $\rho<1$ every time, the error
is at most the update divided by $1-\rho$.  For a slowly converging sweep,
$\rho$ near one, that is many times the tolerance.

\begin{theorem}[T3d, stopping on the update]\label{thm:T3d}
  \lean{Atlas.Schwarz.error_le_update}
  \leanok
  \uses{thm:T7, lem:T3c-sweep}
  Let $V$ be a real normed space and $\norm{G_0(e)}\le\rho\norm{e}$ for all
  $e$, with $\rho<1$.  If $Au^\star=f$ then for every $u\in V$
  \[
    \norm{u-u^\star}\le\frac{\norm{G_f(u)-u}}{1-\rho}.
  \]
\end{theorem}

\begin{proof}
  \uses{thm:T7, lem:T3c-sweep}
  \leanok
  $G_f(u)-G_f(u^\star)=G_0(u-u^\star)$ and $G_f(u^\star)=u^\star$, so $G_f$
  contracts the pair $(u,u^\star)$ by $\rho$; apply Theorem~\ref{thm:T7}.
\end{proof}
